% Appendix A: Full Indicator Definitions
\section{Full Indicator Definitions}\label{app:indicators}

Each indicator measures a specific architectural feature predicted by one of the seven consciousness theories.
This appendix provides the complete definition, scoring methodology, and threshold rationale for all 20 indicators.

%% ─── A.1  GWT ────────────────────────────────────────────────────────────────
\subsection{Global Workspace Theory (GWT)}

\subsubsection{\texttt{global\_broadcast}}
\begin{description}
  \item[Theory:] GWT \cite{baars1988,dehaene2001}
  \item[Weight:] 1.2
  \item[Threshold:] 0.5
  \item[What it measures:] Whether workspace winners are distributed to all registered downstream modules after ignition.
  \item[Scoring:] Ratio of modules that received the broadcast to total registered modules. Score of 1.0 means all modules received broadcast content.
  \item[Threshold rationale:] Set at 0.5 because partial broadcast (reaching at least half of modules) still constitutes meaningful global access.
\end{description}

\subsubsection{\texttt{ignition\_dynamics}}
\begin{description}
  \item[Theory:] GWT \cite{dehaene2001}
  \item[Weight:] 1.0
  \item[Threshold:] 0.4
  \item[What it measures:] Whether candidates undergo nonlinear amplification when activation exceeds the ignition threshold.
  \item[Scoring:] Based on the count and magnitude of ignition events in recent cycles. Higher scores indicate more frequent and stronger ignition.
  \item[Threshold rationale:] Set at 0.4 to allow detection even with limited cycle history.
\end{description}

\subsubsection{\texttt{global\_ignition\_nuanced}}
\begin{description}
  \item[Theory:] GWT
  \item[Weight:] 1.0
  \item[Threshold:] 0.4
  \item[What it measures:] Refined ignition dynamics including sustain duration and decay patterns, beyond the binary ignition detection of the basic indicator.
  \item[Scoring:] Composite of ignition frequency, sustain duration, and decay rate. Higher scores indicate more biologically plausible ignition dynamics.
  \item[Threshold rationale:] Matched to \texttt{ignition\_dynamics} threshold for consistency.
\end{description}

\subsubsection{\texttt{recurrent\_processing}}
\begin{description}
  \item[Theory:] GWT (shared with RPT)
  \item[Weight:] 0.9
  \item[Threshold:] 0.4
  \item[What it measures:] Feedback loops between modules---whether output from later-stage modules feeds back to influence earlier-stage processing.
  \item[Scoring:] Based on detected feedback pathways in the consciousness cycle (e.g., FEP prediction errors modulating GWT salience in subsequent cycles).
  \item[Threshold rationale:] Set at 0.4 as recurrence is present by architectural design in the cross-cycle feedback loop.
\end{description}

%% ─── A.2  AST ────────────────────────────────────────────────────────────────
\subsection{Attention Schema Theory (AST)}

\subsubsection{\texttt{attention\_schema}}
\begin{description}
  \item[Theory:] AST \cite{graziano2013}
  \item[Weight:] 1.1
  \item[Threshold:] 0.5
  \item[What it measures:] Whether the system maintains a self-model of its own attention---not the attention mechanism itself, but a predictive model of where attention is directed and where it will shift.
  \item[Scoring:] Composite of schema completeness (current focus tracked, state classified, history maintained) and prediction accuracy (does the schema correctly predict attention shifts).
  \item[Threshold rationale:] Set at 0.5 because AST requires both schema existence and predictive accuracy.
\end{description}

\subsubsection{\texttt{attention\_control}}
\begin{description}
  \item[Theory:] AST
  \item[Weight:] 1.0
  \item[Threshold:] 0.4
  \item[What it measures:] Whether the system can voluntarily shift attention rather than only responding to bottom-up capture.
  \item[Scoring:] Ratio of voluntary to captured attention shifts. Higher scores indicate more volitional control.
  \item[Threshold rationale:] Set at 0.4 because some voluntary control is expected even with limited operational history.
\end{description}

\subsubsection{\texttt{embodiment}}
\begin{description}
  \item[Theory:] AST (supplementary)
  \item[Weight:] 0.8
  \item[Threshold:] 0.3
  \item[What it measures:] Whether the system models its own boundaries and presence---a sense of being a bounded entity distinct from its environment.
  \item[Scoring:] Based on the presence and coherence of homeostatic drive monitoring, body state tracking (Damasio protoself), and self-world boundary modeling.
  \item[Threshold rationale:] Low threshold (0.3) because embodiment in the standalone framework is simulated rather than genuine.
\end{description}

%% ─── A.3  HOT ────────────────────────────────────────────────────────────────
\subsection{Higher-Order Thought Theory (HOT)}

\subsubsection{\texttt{higher\_order\_representations}}
\begin{description}
  \item[Theory:] HOT \cite{rosenthal2005}
  \item[Weight:] 1.2
  \item[Threshold:] 0.5
  \item[What it measures:] Whether workspace winners receive meta-representations---higher-order thoughts about first-order mental states.
  \item[Scoring:] Ratio of workspace winners that receive at least one HOT to total workspace winners. Score of 1.0 means every winner has a meta-representation.
  \item[Threshold rationale:] Set at 0.5 because HOT theory requires meta-representation for conscious content, but not every processing step may generate one.
\end{description}

\subsubsection{\texttt{metacognition}}
\begin{description}
  \item[Theory:] HOT
  \item[Weight:] 1.1
  \item[Threshold:] 0.4
  \item[What it measures:] Whether the system reflects on its own cognitive processes---monitoring, evaluating, and adjusting its own reasoning.
  \item[Scoring:] Based on introspection depth, self-model accuracy, and the presence of meta-uncertainty (uncertainty about its own confidence levels).
  \item[Threshold rationale:] Set at 0.4 to detect metacognitive activity even at shallow recursion depth.
\end{description}

%% ─── A.4  FEP ────────────────────────────────────────────────────────────────
\subsection{Free Energy Principle (FEP)}

\subsubsection{\texttt{prediction\_error\_minimization}}
\begin{description}
  \item[Theory:] FEP \cite{friston2010}
  \item[Weight:] 1.0
  \item[Threshold:] 0.4
  \item[What it measures:] Whether prediction errors decrease through belief updating---the core operation of the Free Energy Principle.
  \item[Scoring:] Based on the trend of prediction error over recent cycles. Decreasing errors score higher; increasing or stagnant errors score lower.
  \item[Threshold rationale:] Set at 0.4 because the hierarchical predictive processor generates structural prediction errors even from cold start.
\end{description}

\subsubsection{\texttt{hierarchical\_prediction}}
\begin{description}
  \item[Theory:] FEP
  \item[Weight:] 1.0
  \item[Threshold:] 0.4
  \item[What it measures:] Whether prediction operates across multiple hierarchy levels---each level generating predictions about the level below and receiving prediction errors from it.
  \item[Scoring:] Based on the number of active hierarchy levels and the presence of inter-level prediction error flow.
  \item[Threshold rationale:] The 3--4 level hierarchy is structural, so this indicator scores high from initialization.
\end{description}

\subsubsection{\texttt{agency}}
\begin{description}
  \item[Theory:] FEP
  \item[Weight:] 1.0
  \item[Threshold:] 0.5
  \item[What it measures:] Whether the system selects actions to fulfill predictions---active inference rather than passive observation.
  \item[Scoring:] Based on action selection frequency, policy quality, and the presence of goal-directed behavior in the active inference cycle.
  \item[Threshold rationale:] Set at 0.5 because genuine agency requires demonstrable action selection, not merely prediction.
\end{description}

\subsubsection{\texttt{sparse\_smooth\_coding}}
\begin{description}
  \item[Theory:] FEP (supplementary)
  \item[Weight:] 0.8
  \item[Threshold:] 0.3
  \item[What it measures:] Whether internal representations are sparse (few active units) and smooth (similar inputs produce similar representations).
  \item[Scoring:] Composite of sparsity (fraction of near-zero activations) and smoothness (gradient continuity of representations with respect to input).
  \item[Threshold rationale:] Low threshold because the FEP implementation structurally produces sparse representations.
\end{description}

%% ─── A.5  IIT ────────────────────────────────────────────────────────────────
\subsection{Integrated Information Theory (IIT)}

\subsubsection{\texttt{integrated\_information}}
\begin{description}
  \item[Theory:] IIT \cite{tononi2008}
  \item[Weight:] 1.3
  \item[Threshold:] 0.3
  \item[What it measures:] Whether the system generates integrated information ($\Phi > 0$)---the whole produces more information than the sum of its parts.
  \item[Scoring:] Normalized $\Phi$ value from the IIT approximation (exact computation via PyPhi is unavailable in Python 3.10+). Score is $\Phi$ mapped to $[0, 1]$.
  \item[Threshold rationale:] Low threshold (0.3) because the approximation produces conservative estimates and any non-trivial $\Phi$ is significant.
\end{description}

\subsubsection{\texttt{irreducibility}}
\begin{description}
  \item[Theory:] IIT
  \item[Weight:] 1.0
  \item[Threshold:] 0.4
  \item[What it measures:] Whether the system cannot be decomposed into independent parts without information loss---a core axiom of IIT.
  \item[Scoring:] Based on comparing the information generated by the full system to the information generated by its partitions. Higher scores indicate stronger irreducibility.
  \item[Threshold rationale:] Set at 0.4 because meaningful irreducibility requires substantial inter-module information flow.
\end{description}

%% ─── A.6  RPT ────────────────────────────────────────────────────────────────
\subsection{Recurrent Processing Theory (RPT)}

\subsubsection{\texttt{local\_recurrence}}
\begin{description}
  \item[Theory:] RPT \cite{lamme2006}
  \item[Weight:] 0.8
  \item[Threshold:] 0.3
  \item[What it measures:] Feedback within individual neural substrates---recurrent connections within the SNN and LSM that allow local re-processing of signals.
  \item[Scoring:] Based on recurrent weight strength in the SNN (STDP-learned feedback connections) and reservoir dynamics in the LSM (spectral radius as proxy for recurrence).
  \item[Threshold rationale:] Low threshold because even structurally recurrent networks score non-zero.
\end{description}

\subsubsection{\texttt{algorithmic\_recurrence}}
\begin{description}
  \item[Theory:] RPT
  \item[Weight:] 0.9
  \item[Threshold:] 0.3
  \item[What it measures:] Computational recurrence measured in the substrate dynamics---not just structural feedback connections, but active recurrent processing.
  \item[Scoring:] Based on spike re-entry rates (SNN), reservoir state trajectory analysis (LSM), and the classification of recurrence type (superficial vs.\ deep).
  \item[Threshold rationale:] Set at 0.3 but requires active neural substrates to score above baseline; scores low without them.
\end{description}

%% ─── A.7  BLT ────────────────────────────────────────────────────────────────
\subsection{Beautiful Loop Theory (BLT)}

\subsubsection{\texttt{bayesian\_binding\_quality}}
\begin{description}
  \item[Theory:] BLT \cite{laukkonen2025}
  \item[Weight:] 1.0
  \item[Threshold:] 0.3
  \item[What it measures:] Whether separate inferences bind into a coherent unified percept---measured through mutual information between workspace winners.
  \item[Scoring:] Based on the binding quality score from the \texttt{BayesianBinding} module: the total mutual information in the selected coherent subset, normalized by the number of bound inferences.
  \item[Threshold rationale:] Low threshold because binding quality depends on having multiple workspace winners with correlated beliefs, which requires accumulated history.
\end{description}

\subsubsection{\texttt{epistemic\_depth}}
\begin{description}
  \item[Theory:] BLT
  \item[Weight:] 1.1
  \item[Threshold:] 0.3
  \item[What it measures:] The depth of recursive self-reference---how many levels of self-modeling the system exhibits (0 = none, 1 = self-model, 2 = meta-predictions, 3 = HOTs about self, 4+ = higher-order reflection).
  \item[Scoring:] Depth level divided by 4, capped at 1.0. A depth of 2 scores 0.5; a depth of 4 scores 1.0.
  \item[Threshold rationale:] Set at 0.3 (equivalent to depth $\geq$ 1.2), requiring at least a basic self-model to pass.
\end{description}

\subsubsection{\texttt{genuine\_implementation}}
\begin{description}
  \item[Theory:] BLT
  \item[Weight:] 1.2
  \item[Threshold:] 0.4
  \item[What it measures:] Anti-mimicry check---whether the Beautiful Loop components are producing results consistent with genuine recursive processing rather than merely reporting high numerical scores.
  \item[Scoring:] Requires three conditions to score above threshold:
    (1)~field-evidencing detected (the system recursively evidences its own predictions),
    (2)~loop quality above 0.4, and
    (3)~a strange loop detected by the epistemic depth tracker.
    Meeting all three scores 1.0; partial scores proportional to conditions met.
  \item[Threshold rationale:] Set at 0.4 to require at least two of the three genuine-implementation conditions. This is the assessment's primary defense against score inflation.
\end{description}
